\documentclass[12pt,a4paper]{article}

\usepackage[utf8]{inputenc}
\usepackage[T1]{fontenc}
\usepackage[french]{babel}
\usepackage{lmodern}
\usepackage{amsmath,amssymb}
\usepackage{geometry}
\usepackage{microtype}
\usepackage{hyperref}
\usepackage{enumitem}
\geometry{top=2.5cm,bottom=2.5cm,left=2.5cm,right=2.5cm}
\hypersetup{hidelinks,pdftitle={Documentation Replay Buffer HCANN}}

\title{Documentation du Buffer Épisodique}
\author{Projet HCANN-Agent}
\date{\today}

\begin{document}
\maketitle

\section{Module Tampon (\texttt{replay\_buffer.py})}

\subsection{Objectif du module}
Le module fournit un tampon FIFO pour stocker temporairement les épisodes avant consolidation lente cortex$\rightarrow$hippocampe.

\begin{itemize}
    \item \textbf{EpisodicBuffer}: buffer principal avec capacité bornée;
    \item \textbf{ExperienceReplayBuffer}: alias de compatibilité avec du code existant.
\end{itemize}

\subsection{Structure d'un épisode}
Chaque entrée stockée contient:
\begin{itemize}
    \item \texttt{sem}: embedding sémantique (tenseur PyTorch);
    \item \texttt{hpc}: état hippocampique (tenseur PyTorch);
    \item \texttt{meta}: métadonnées (dictionnaire);
    \item \texttt{ts}: timestamp (float Unix).
\end{itemize}

\subsection{Politique mémoire}
Le buffer s'appuie sur \texttt{collections.deque(maxlen=capacity)}:
\begin{equation}
|\mathcal{B}| \le C,
\end{equation}
où $C$ est la capacité.  
Lorsqu'un nouvel épisode arrive et que $|\mathcal{B}|=C$, le plus ancien est évincé automatiquement (FIFO).

\subsection{Vue théorique (mémoire de rejeu)}
Le buffer réalise une fenêtre temporelle glissante:
\begin{equation}
\mathcal{B}_t = \{x_{t-C+1}, \dots, x_t\},
\end{equation}
avec un épisode:
\begin{equation}
x_t=(\mathbf{s}_t,\mathbf{h}_t,m_t,\tau_t).
\end{equation}
Ici, $\mathbf{s}_t$ est l'état sémantique, $\mathbf{h}_t$ l'état hippocampique, $m_t$ les métadonnées et $\tau_t$ le temps.

L'entraînement lent approxime une espérance empirique:
\begin{equation}
\mathcal{L}_{\text{slow}}
= \mathbb{E}_{x\sim \mathcal{D}_{recent}}[\ell(x)]
\approx \frac{1}{n}\sum_{x_i \sim \mathcal{B}_t}\ell(x_i).
\end{equation}

\subsection{1) Ajout d'épisode}
\paragraph{\texttt{push(sem\_embed, hpc\_state, metadata=None)}}
\begin{itemize}[leftmargin=*]
    \item détache les tenseurs du graphe de gradients (\texttt{detach()});
    \item les copie sur CPU (\texttt{cpu()}) pour alléger la VRAM;
    \item ajoute une entrée timestampée dans le buffer.
\end{itemize}

\subsection{2) Échantillonnage pour consolidation}
\paragraph{\texttt{sample(n=32)}}
\begin{itemize}[leftmargin=*]
    \item si $|\mathcal{B}| < n$: retourne \texttt{None};
    \item sinon: tirage aléatoire uniforme sans remise de $n$ épisodes;
    \item sortie sous forme batch:
\end{itemize}
\begin{align}
\texttt{sem} &\in \mathbb{R}^{n \times d_s},\\
\texttt{hpc} &\in \mathbb{R}^{n \times d_h}.
\end{align}
Les champs \texttt{meta} et \texttt{ts} sont retournés en listes parallèles.

Si la taille courante est $M=|\mathcal{B}|$ et $n\le M$, la probabilité d'inclusion d'un épisode sous tirage uniforme sans remise est:
\begin{equation}
\mathbb{P}(x_i \in \text{batch}) = \frac{n}{M}.
\end{equation}

\subsection{3) Accès récent et maintenance}
\begin{itemize}
    \item \texttt{get\_recent(n)}: retourne les $n$ épisodes les plus récents;
    \item \texttt{clear()}: vide complètement le buffer;
    \item \texttt{\_\_len\_\_()}: taille courante.
\end{itemize}

\subsection{Correspondance implémentation}
\begin{itemize}
    \item \texttt{\_\_init\_\_}: création de \texttt{deque(maxlen=capacity)};
    \item \texttt{push}: sérialisation légère (\texttt{detach/cpu}) + insertion;
    \item \texttt{sample}: \texttt{random.sample} + \texttt{torch.stack};
    \item \texttt{ExperienceReplayBuffer}: héritage direct sans surcharge.
\end{itemize}

\subsection{Rendu computationnel (coûts)}
\begin{itemize}
    \item \texttt{push}: $O(1)$ amorti (ajout et éventuelle éviction FIFO);
    \item \texttt{sample}: $O(n)$ pour le tirage + $O(n d_s + n d_h)$ pour la construction tensorielle;
    \item \texttt{get\_recent(n)}: $O(n)$.
\end{itemize}

Correspondance théorie$\leftrightarrow$implémentation:
\begin{itemize}
    \item mémoire glissante bornée $\rightarrow$ \texttt{deque(maxlen=capacity)};
    \item rejeu stochastique $\rightarrow$ \texttt{random.sample};
    \item mini-batch pour consolidation $\rightarrow$ \texttt{torch.stack} des composantes \texttt{sem/hpc}.
\end{itemize}

\subsection{Interface détaillée}

\paragraph{\texttt{EpisodicBuffer(capacity=2000)}}
\begin{itemize}[leftmargin=*]
    \item \textbf{Entrée}: capacité maximale.
    \item \textbf{Effet}: initialise un buffer borné FIFO.
\end{itemize}

\paragraph{\texttt{push(sem\_embed, hpc\_state, metadata=None)}}
\begin{itemize}[leftmargin=*]
    \item \textbf{Entrées}: deux tenseurs + métadonnées optionnelles.
    \item \textbf{Effet}: ajoute un épisode prêt pour le replay.
\end{itemize}

\paragraph{\texttt{sample(n=32)}}
\begin{itemize}[leftmargin=*]
    \item \textbf{Entrée}: taille batch souhaitée.
    \item \textbf{Sortie}: dictionnaire de batch ou \texttt{None}.
\end{itemize}

\paragraph{\texttt{get\_recent(n=10)}}
\begin{itemize}[leftmargin=*]
    \item \textbf{Sortie}: liste des derniers épisodes (ordre chronologique).
\end{itemize}

\paragraph{\texttt{clear()}}
\begin{itemize}[leftmargin=*]
    \item \textbf{Effet}: supprime tout le contenu du tampon.
\end{itemize}

\paragraph{\texttt{ExperienceReplayBuffer}}
\begin{itemize}[leftmargin=*]
    \item \textbf{Rôle}: compatibilité descendante; API identique à \texttt{EpisodicBuffer}.
\end{itemize}

\subsection{Intégration dans la boucle d'apprentissage}
Cycle typique:
\begin{enumerate}[leftmargin=*]
    \item encodage rapide d'un épisode $(sem, hpc)$;
    \item insertion via \texttt{push};
    \item tous les \texttt{consolidation\_freq} pas, tirage via \texttt{sample};
    \item optimisation de la perte de consolidation lente sur le batch.
\end{enumerate}

\end{document}
